\documentclass[10pt,a4paper]{amsart}
\usepackage[margin=19mm]{geometry}
\usepackage{amsmath,amssymb,mathtools}
\usepackage[scr=rsfs,cal=euler]{mathalfa}
\usepackage{xfrac}
\usepackage[unicode,hidelinks,bookmarks=false]{hyperref}
\newcommand{\ie}{i.e.\ }
\newcommand{\cf}{cf.\ }
\newcommand{\resp}{respectively\ }
\newcommand{\Subsection}{\subsection}
\providecommand{\nobreakdash}{\nobreak\hbox{-}}
\setlength{\emergencystretch}{2em}
\multlinegap0pt
\usepackage{bm}
\usepackage{bbm}
\usepackage{dsfont}
\usepackage{xspace}
\usepackage{cancel}
\usepackage{mathtools}
\usepackage{amsthm}
\makeatletter
\@ifundefined{definition}{\newtheorem{definition}{Definition}}{}
\@ifundefined{lemma}{\newtheorem{lemma}[definition]{Lemma}}{}
\makeatother
\makeatletter
\expandafter\ifx\csname OrangeEnv\endcsname\relax
\newenvironment{OrangeEnv}{\color{orange}}{}
\fi
\DeclareMathOperator{\arcosh}{arcosh}
\expandafter\ifx\csname inlineenum\endcsname\relax
\newenvironment{inlineenum}
{\setcounter{inlineenum}{0}%
\renewcommand{\item}{\refstepcounter{inlineenum}{(\theinlineenum)\hspace{\thelabelsep}}}
}
\fi
\newcommand{\lm}[1]{\mathbb{L}^2(#1)}
\newcommand{\ma}{\ensuremath{\measuredangle}}
\makeatother
\title[Total curvature and length estimates for timelike curves in ]{Total curvature and length estimates for timelike curves in Lorentzian length spaces}
\author{Darius Er\"os, Felix Rott, Zhe-Feng Xu}
\date{Source: arXiv:2606.17834v1 (CC BY 4.0)}
\begin{document}
\maketitle

\noindent\emph{Source and licence.} Total curvature and length estimates for timelike curves in Lorentzian length spaces. Version arXiv:2606.17834v1 (CC BY 4.0), distributed under the Creative Commons Attribution 4.0 International licence, as recorded in the arXiv licence metadata for this exact version and shown on the abstract page https://arxiv.org/abs/2606.17834v1. Full rights notice: licenses/arxiv-2606.17834-CC-BY-4.0.txt.\par\smallskip

\noindent\textit{Editorial scope.} Counted unit: the Lemma whose source label is \texttt{lem:positive\_tri\_segment\_estimate} in the article (source0.tex lines 1318--1327, source label \texttt{lem:positive\_tri\_segment\_estimate}), delivered together with the article's own complete original proof of it (source0.tex lines 1329--1497, one \texttt{proof} environment of 169 lines). No mathematics is written, repaired or re-derived: statement and proof are the article's own text line for line. Required context retained uncounted: lem: bi-segment and rigidity (lines 1060--1088). Every definition, display and label of those lines is retained unchanged. Omitted from this package and not redistributed: 1--1059, 1089--1317, 1328--1328, 1498--1822. Nothing inside a retained excerpt is omitted or altered: every retained body is delivered exactly as the article's lines are written, so no omission receipt of any kind accompanies this package. Citations: the retained excerpts carry no citation command at all, so no bibliography entry of the article is transcribed and no reference is redistributed. Reference adaptation: none was needed and none was made; every reference inside the delivered unit resolves inside it, so no unresolved reference or citation is produced. Printed slips kept literal: no printed word, symbol, hypothesis or number of the counted statement or of its proof was altered. Layout and class: the article's own class is \texttt{article} with packages including inputenc, geometry, amssymb, amsmath, bm, mathtools, amsthm, babel, lmodern, microtype, euscript, zref-clever, hyperref, xcolor; the wrapper is the lane's pinned \texttt{amsart} editorial wrapper at 10pt on A4, which restates the theorem-like environments the retained text needs and adds the article's own macro definitions that the retained text needs (\texttt{\textbackslash arcosh}, \texttt{\textbackslash lm}, \texttt{\textbackslash ma}), with extra wrapper packages (bm, bbm, dsfont, xspace, cancel, mathtools). The delivered numbering is the unit's own. Assets: no retained span contains a figure environment, a graphics-inclusion command, a PDF-page inclusion, a table or a TikZ picture; no image file is copied into this package, the package declares no asset dependency and no image file is redistributed with it. Licence: version arXiv:2606.17834v1 of this article is distributed under the Creative Commons Attribution 4.0 International licence, as recorded in the arXiv licence metadata for this exact version and shown on the abstract page https://arxiv.org/abs/2606.17834v1. A full rights notice is delivered at licenses/arxiv-2606.17834-CC-BY-4.0.txt. Characters: every retained excerpt is pure ASCII, so the delivered engine, XeLaTeX with its default Latin Modern text font, reports no missing character.
\begin{lemma}[Length of bi-segments and rigidity]
\label{lem: bi-segment and rigidity}
Let $p\ll q\ll w$ be three timelike related points in $\lm{K}$. 
Let $\kappa=\ma_{q}(p,w)$ be the total curvature of the corresponding bi-segment, and let $r=\ell(p,w)<D_K$. 
Assume that $\ell(p,q)+\ell(q,w)=l_0$ is fixed. 
Then
\begin{equation*}
\label{total bound}
\kappa\geq \bm{\kappa}(K,r,l_0):=
\begin{cases}
2\arcosh\left(\dfrac{\sin(\sqrt{K}r/2)}{\sin(\sqrt{K}l_0/2)}\right), 
& \mathrm{if}\,\,K>0,\\[1.2ex]
2\arcosh\left(\dfrac{r}{l_0}\right), 
& \mathrm{if}\,\,K=0,\\[1.2ex]
2\arcosh\left(\dfrac{\sinh(\sqrt{-K}r/2)}{\sinh(\sqrt{-K}l_0/2)}\right), 
& \mathrm{if}\,\,K<0.
\end{cases}
\end{equation*}
Equality holds if and only if $\ell(p,q)=\ell(q,w)=l_0/2$ when $\kappa>0$, and if and only if $p,q,w$ lie on a geodesic in the case $\kappa=0$. Moreover, if $\kappa_0\in[0,+\infty)$ denotes the total curvature of the corresponding isosceles bi-segment, then its length is given by
\begin{equation}
\label{length bound}
\mathcal{L}(K,r,\kappa_0):=
\begin{cases}
\dfrac{2}{\sqrt{K}}\arcsin\left(\dfrac{\sin(\sqrt{K}r/2)}{\cosh(\kappa_0/2)}\right), & \mathrm{if}\,\,K>0,\\[1.2ex]
\dfrac{r}{\cosh(\kappa_0/2)}, & \mathrm{if}\,\,K=0,\\[1.2ex]
\dfrac{2}{\sqrt{-K}}\operatorname{arsinh}\left(\dfrac{\sinh(\sqrt{-K}r/2)}{\cosh(\kappa_0/2)}\right), & \mathrm{if}\,\,K<0.
\end{cases}
\end{equation}
\end{lemma}

\begin{lemma}[Direct estimate for $K > 0$]
\label{lem:positive_tri_segment_estimate}
Let $K>0$ and let $p_0\ll p_1\ll p_2\ll p_3$ be a convex timelike poly-segment consisting of three segments in $\lm{K}$ with total length $l$ and $ r:=\ell(p_0,p_3) < D_K$.\footnote{Here, convex means that the quadrilateral generated by these four vertices is convex.} Assume $\theta=\ma_{p_1}(p_0,p_2)>0,
\varphi=\ma_{p_2}(p_1,p_3)>0,$ and the total curvature $\Theta=\theta+\varphi$.
Then there exists an isosceles timelike bi-segment $B$ in $\lm{K}$ with the same endpoints
$p_0$ and $p_3$, total length $l$, and
\[
\kappa^*(B)< \Theta .
\]
\end{lemma}

\begin{proof}
Set
\begin{equation*}
\begin{aligned}
&a:=\ell(p_0,p_1), &&b:=\ell(p_1,p_2), &&c\;:=\ell(p_2,p_3), 
&d:=\ell(p_0,p_2), &&\psi:=\ma_{p_2}(p_0,p_1).
\end{aligned}
\end{equation*}
We claim
\begin{equation*}
\sin\left(\frac{\sqrt K\,r}{2}\right)
< \cosh\left(\frac{\Theta}{2}\right)
\sin\left(\frac{\sqrt K\,l}{2}\right).
\end{equation*}
% Equivalently,
% \[
% 2\,\arcosh
% \left(
% \frac{\sin(\sqrt K\,r/2)}
%      {\sin(\sqrt K\,l/2)}
% \right)
% \le \Theta .
% \]

By rescaling, it suffices to prove the claim for $K=1$, in which case $D_K=\pi$.  
Moreover, all side lengths appearing below are strictly smaller than $\pi$, and their $\sin$-values are positive.
By the Law of Cosines applied to the triangle $\Delta(p_0,p_1,p_2)$ at the vertices $p_1$ and $p_2$, respectively, we obtain 
\begin{equation}\label{eq add 1}
\cos d =
\cos a\,\cos b-\cosh\theta\,\sin a\, \sin b,\qquad \cos a =
\cos b\,\cos d+\cosh\psi\,\sin b\, \sin d.
\end{equation}
Combining these two identities yields
\begin{equation}\label{eq add 2}
\cosh\psi\,\sin d
=
\cos a\,\sin b+\cosh\theta\,\sin a\,\cos b.
\end{equation}
Moreover, solving for the respective $\cosh$-terms, squaring them and using $\cosh^2-\sinh^2=1$, we get
\begin{equation}\label{eq add 3}
\sinh\psi\,\sin d
=
\sinh\theta\,\sin a.
\end{equation}
Applying the Law of Cosines in $\Delta(p_0,p_2,p_3)$ at $p_2$ and using $\ma_{p_2}(p_0,p_3) = \ma_{p_2}(p_0,p_1) + \ma_{p_2}(p_1,p_3) = \psi + \varphi$, which holds by the assumed convexity, we obtain
\begin{equation*}
\cos r
=
\cos d\cos c-\cosh(\psi+\varphi)\sin d\sin c.    
\end{equation*}
Combining this with \eqref{eq add 1}, \eqref{eq add 2} and \eqref{eq add 3}, we get
\begin{equation*}
\begin{aligned}
\cos r
=\;&
\bigl(\cos a\cos b-\cosh\theta\,\sin a\sin b\bigr)\cos c\\
&-\cosh\varphi\,
\bigl(\cos a\sin b+\cosh\theta\,\sin a\cos b\bigr)\sin c\\
&-\sinh\theta\sinh\varphi\,\sin a\sin c.
\end{aligned}
\end{equation*}
Then, using the identity
\begin{equation*}
\cos(a+b+c)
=
\cos a\,\cos b\,\cos c
-\sin a\,\sin b\,\cos c
-\cos a\,\sin b\,\sin c
-\sin a\,\cos b\,\sin c,
\end{equation*}
we obtain
\begin{equation*}
\begin{aligned}
\cos r
=\;&
\cos(a+b+c)
-(\cosh\theta-1)\sin a\,\sin(b+c)\\
&-(\cosh\varphi-1)\sin c\,\sin(a+b)\\
&-\Big((\cosh\theta-1)(\cosh\varphi-1)\cos b
+\sinh\theta\sinh\varphi\Big)\sin a\,\sin c.
\end{aligned}
\end{equation*}
We now estimate the right-hand side. Since $a+b+c=l<\pi,$
we have
\begin{equation*}
\sin a\, \sin(b+c)\leq \sin^2\frac{l}{2},
\qquad
\sin c\, \sin(a+b)\leq \sin^2\frac{l}{2},\qquad \sin a\, \sin c< \sin^2\frac{l}{2}.
\end{equation*}
Set
\begin{equation*}
C:=(\cosh\theta-1)(\cosh\varphi-1)\cos b
+\sinh\theta\sinh\varphi.
\end{equation*}
Then $C > 0$. Indeed,
\begin{equation*}
\begin{aligned}
C&>
-(\cosh\theta-1)(\cosh\varphi-1)
+\sinh\theta\sinh\varphi\\
&=
4\sinh\frac\theta2\sinh\frac{\varphi}{2}
\cosh\frac{\theta-\varphi}{2}
>0.
\end{aligned}
\end{equation*}
Moreover, 
\begin{equation*}
C<(\cosh\theta-1)(\cosh\varphi-1)
+\sinh\theta\sinh\varphi.
\end{equation*}
Therefore,
\begin{equation*}
\begin{aligned}
\cos r
&>
\cos l-\sin^2\frac{l}{2}
\Big[(\cosh\theta-1)+(\cosh\varphi-1)
+(\cosh\theta-1)(\cosh\varphi-1)
+\sinh\theta\sinh\varphi\Big]\\
&=\cos l-(\cosh\Theta-1)\sin^2\frac{l}{2}\\
&=1-2\cosh^2\frac{\Theta}{2}\,\sin^2\frac{l}{2}.
\end{aligned}
\end{equation*}
Equivalently,
\begin{equation*}
1-2\sin^2\frac{r}{2}
>
1-2\cosh^2\frac{\Theta}{2}\,\sin^2\frac{l}{2},
\end{equation*}
and hence
\begin{equation*}
\sin\frac{r}{2}<
\cosh\frac{\Theta}{2}\,\sin\frac{l}{2}.
\end{equation*}
Reintroducing the factor $\sqrt K$, we get
\begin{equation*}
\sin\left(\frac{\sqrt K\,r}{2}\right)<\cosh\left(\frac{\Theta}{2}\right)
\sin\left(\frac{\sqrt K\,l}{2}\right),
\end{equation*}
by the reverse triangle inequality, we know that $l\leq r<D_K$, so the function $t\mapsto \sin(\sqrt K\,t/2)$ is increasing on
$[0,D_K]$, and therefore
\begin{equation*}
1\leq\frac{\sin(\sqrt K\,r/2)}{\sin(\sqrt K\,l/2)}.
\end{equation*}
Defining
\begin{equation*}
\kappa_0:=2\,\arcosh\left(
\frac{\sin(\sqrt K\,r/2)}{\sin(\sqrt K\,l/2)}\right),
\end{equation*}
we get $\kappa_0< \Theta$. 
Now, by Lemma \ref{lem: bi-segment and rigidity}, $\kappa_0$ is precisely the total curvature of an isosceles bi-segment where the endpoints are a distance $r$ apart, and the sides are each of length $l/2$. \qedhere

% \fel{Rest is not needed?}

% Finally, by law of cosines in $L^2(K)$, an isosceles timelike bi-segment with side lengths $l/2$ and vertex angle $\kappa_0$ has endpoint separation $r$, we have
% \[
% \cos(\sqrt K\,r)
% =
% \cos^2\!\left(\frac{\sqrt K\,l}{2}\right)
% -
% \cosh(\kappa_0)
% \sin^2\!\left(\frac{\sqrt K\,l}{2}\right).
% \]
% Thus such an isosceles timelike bi-segment $B$ exists, has total length $l$, and satisfies
% \[
% \kappa^*(B)=\kappa_0\le\Theta .
% \]
\end{proof}

\end{document}
